\chapter{Conclusions}

This thesis has presented a comprehensive study on the detection of phishing emails through linguistic patterns and sentiment analysis. The work encompassed dataset creation and annotation, the training and evaluation of transformer-based classification models, and the development of a deployable web application with a companion browser extension.

\section{Summary of Contributions}

The dataset creation process leveraged advanced \acs{LLM} techniques to generate a diverse and representative corpus of phishing emails across multiple emotional categories. The annotation framework was designed to optimize inter-annotator agreement and minimise bias, resulting in high-quality labelled data. A principled consolidation of the original fifteen-category emotional taxonomy into eight semantically coherent classes was carried out, informed by co-occurrence analysis of the annotated corpus.

Two distinct classification tasks were addressed. For phishing detection, fine-tuned \acs{RoBERTa}-large and \acs{ELECTRA}-large binary classifiers achieved F1-scores of 0.9150 and 0.9055, respectively, on the held-out test set. A three-tier ensemble inference scheme,  classifying emails as \textit{ham}, \textit{maybe spam}, or \textit{spam}, was adopted to surface borderline cases without forcing a hard binary decision, reflecting the asymmetric cost of false negatives in a security context.

For emotion detection, the same two architectures were fine-tuned as multi-label classifiers over the eight consolidated emotional classes. Per-class threshold tuning on the validation set yielded macro-F1 scores of 0.7041 (\acs{RoBERTa}-large) and 0.6967 (\acs{ELECTRA}-large) on the test set, demonstrating that transformer models can reliably identify the psychological manipulation tactics present in phishing emails.

The trained models were integrated into a full-stack web application, \textit{Gone Phishing}, consisting of a FastAPI backend hosted on Hugging Face Spaces and a static frontend that accepts \texttt{.eml} file uploads. A companion Manifest V3 Chrome extension extends the same analysis capability to Gmail and Outlook Web, allowing users to scan open emails directly from their inbox without leaving the browser.

\section{Limitations}

Several limitations should be acknowledged. The phishing email corpus was generated synthetically using \acs{LLM} prompting, which may not fully capture the diversity and adversarial sophistication of real-world phishing campaigns. The annotation process was constrained by time and annotator availability, with only a subset of the full corpus receiving manual validation. Emotion detection performance, while promising, remains below the level achieved for binary spam classification, partly due to the inherent subjectivity of multi-label emotional annotation and the class imbalance present even after taxonomy consolidation.

Furthermore, there is a serious risk of methodological circularity in the dataset creation process. The 10\,000 phishing emails were generated by an \acs{LLM} using prompts that explicitly encoded emotional categories, and the same generated emails were then used to train the emotion classifier. In this setting, the model may learn to recover the label implied by the prompt-generation process rather than detect emotion in naturally occurring phishing emails. This can inflate performance estimates on the held-out test set, since the test data is drawn from the same synthetic distribution. Future work should therefore include evaluation on external datasets of real-world phishing emails, ideally with human annotations, to assess whether the model generalises beyond prompt-induced patterns.

\section{Future Work}

Future work should prioritise breaking the circularity introduced by prompt-driven synthetic data generation. This requires collecting real-world phishing examples, ideally from sources such as the \acs{APWG} eCrime Exchange, and annotating them independently so that emotion labels are not implied by the generation prompt. A broader set of annotators would also improve inter-annotator agreement and label coverage. Further investigation into adversarial robustness is warranted, particularly as \acs{AI}-generated phishing content becomes more prevalent. On the application side, extending the Chrome extension to support additional email clients and languages would increase the practical reach of the system.